\section{Velocity Direction}

Choosing an appropriate \emph{direction} for the initial velocity is just as important as choosing its magnitude. If the velocity were directed radially (i.e., along the line joining the star and planet), the planet would simply fall toward or away from the star. For an orbit to form, the velocity must have a component perpendicular to the radial direction.

In two dimensions, the natural choice is to set the initial velocity entirely perpendicular to the radius vector. This gives the planet the best chance of entering a stable, closed orbit.

\subsection{The Radius Vector}

Using the initial positions from the Problem Setup, the displacement vector from the star to the planet is:

\begin{equation}
\vec{r} = \vec{r}_p - \vec{r}_\star = (100 - 0)\,\hat{\mathbf{i}} + (75 - 0)\,\hat{\mathbf{j}} = 100\,\hat{\mathbf{i}} + 75\,\hat{\mathbf{j}}
\end{equation}

with magnitude $r = |\vec{r}| = 125$, as computed in the Find Distance section.

\subsection{Constructing the Perpendicular Vector}

In two dimensions, a vector perpendicular to $(x, y)$ is obtained by rotating it $90°$ counter-clockwise, yielding $(-y, x)$. Applying this to $\vec{r}$:

\begin{equation}
\vec{r}_\perp = -75\,\hat{\mathbf{i}} + 100\,\hat{\mathbf{j}}
\end{equation}

One can verify perpendicularity by checking the dot product:

\begin{equation}
\vec{r} \cdot \vec{r}_\perp = (100)(-75) + (75)(100) = -7500 + 7500 = 0 \checkmark
\end{equation}

\subsection{Computing the Unit Perpendicular Vector}

To obtain a unit vector in the perpendicular direction, divide $\vec{r}_\perp$ by its magnitude (which equals $r$, since rotation preserves length):

\begin{equation}
\hat{r}_\perp = \frac{\vec{r}_\perp}{|\vec{r}_\perp|} = \frac{-75\,\hat{\mathbf{i}} + 100\,\hat{\mathbf{j}}}{125} = -0.6\,\hat{\mathbf{i}} + 0.8\,\hat{\mathbf{j}}
\end{equation}

This is implemented in the simulation code via the \texttt{Vec2::perpendicular()} method combined with division by the scalar distance.

\subsection{Assigning the Initial Velocity Vector}

The initial velocity vector is then the chosen scalar speed $v_0$ multiplied by the unit perpendicular:

\begin{equation}
\vec{v}_0 = v_0 \cdot \hat{r}_\perp
\end{equation}

Substituting $v_0 = 2.8$ for the worked example:

\begin{align}
\vec{v}_0 &= 2.8 \times \left(-0.6\,\hat{\mathbf{i}} + 0.8\,\hat{\mathbf{j}}\right) \\
          &= -1.68\,\hat{\mathbf{i}} + 2.24\,\hat{\mathbf{j}}
\end{align}

The planet therefore begins moving up and to the left, perpendicular to the line joining it to the star. Under the star's gravitational pull, this initial velocity is gradually curved into the orbital trajectory that the Velocity Verlet integrator tracks over subsequent timesteps.